% Options for packages loaded elsewhere
\PassOptionsToPackage{unicode}{hyperref}
\PassOptionsToPackage{hyphens}{url}
\documentclass[
]{article}
\usepackage{xcolor}
\usepackage{amsmath,amssymb}
\setcounter{secnumdepth}{-\maxdimen} % remove section numbering
\usepackage{iftex}
\ifPDFTeX
  \usepackage[T1]{fontenc}
  \usepackage[utf8]{inputenc}
  \usepackage{textcomp} % provide euro and other symbols
\else % if luatex or xetex
  \usepackage{unicode-math} % this also loads fontspec
  \defaultfontfeatures{Scale=MatchLowercase}
  \defaultfontfeatures[\rmfamily]{Ligatures=TeX,Scale=1}
\fi
\usepackage{lmodern}
\ifPDFTeX\else
  % xetex/luatex font selection
\fi
% Use upquote if available, for straight quotes in verbatim environments
\IfFileExists{upquote.sty}{\usepackage{upquote}}{}
\IfFileExists{microtype.sty}{% use microtype if available
  \usepackage[]{microtype}
  \UseMicrotypeSet[protrusion]{basicmath} % disable protrusion for tt fonts
}{}
\makeatletter
\@ifundefined{KOMAClassName}{% if non-KOMA class
  \IfFileExists{parskip.sty}{%
    \usepackage{parskip}
  }{% else
    \setlength{\parindent}{0pt}
    \setlength{\parskip}{6pt plus 2pt minus 1pt}}
}{% if KOMA class
  \KOMAoptions{parskip=half}}
\makeatother
\usepackage{color}
\usepackage{fancyvrb}
\newcommand{\VerbBar}{|}
\newcommand{\VERB}{\Verb[commandchars=\\\{\}]}
\DefineVerbatimEnvironment{Highlighting}{Verbatim}{commandchars=\\\{\}}
% Add ',fontsize=\small' for more characters per line
\newenvironment{Shaded}{}{}
\newcommand{\AlertTok}[1]{\textcolor[rgb]{1.00,0.00,0.00}{\textbf{#1}}}
\newcommand{\AnnotationTok}[1]{\textcolor[rgb]{0.38,0.63,0.69}{\textbf{\textit{#1}}}}
\newcommand{\AttributeTok}[1]{\textcolor[rgb]{0.49,0.56,0.16}{#1}}
\newcommand{\BaseNTok}[1]{\textcolor[rgb]{0.25,0.63,0.44}{#1}}
\newcommand{\BuiltInTok}[1]{\textcolor[rgb]{0.00,0.50,0.00}{#1}}
\newcommand{\CharTok}[1]{\textcolor[rgb]{0.25,0.44,0.63}{#1}}
\newcommand{\CommentTok}[1]{\textcolor[rgb]{0.38,0.63,0.69}{\textit{#1}}}
\newcommand{\CommentVarTok}[1]{\textcolor[rgb]{0.38,0.63,0.69}{\textbf{\textit{#1}}}}
\newcommand{\ConstantTok}[1]{\textcolor[rgb]{0.53,0.00,0.00}{#1}}
\newcommand{\ControlFlowTok}[1]{\textcolor[rgb]{0.00,0.44,0.13}{\textbf{#1}}}
\newcommand{\DataTypeTok}[1]{\textcolor[rgb]{0.56,0.13,0.00}{#1}}
\newcommand{\DecValTok}[1]{\textcolor[rgb]{0.25,0.63,0.44}{#1}}
\newcommand{\DocumentationTok}[1]{\textcolor[rgb]{0.73,0.13,0.13}{\textit{#1}}}
\newcommand{\ErrorTok}[1]{\textcolor[rgb]{1.00,0.00,0.00}{\textbf{#1}}}
\newcommand{\ExtensionTok}[1]{#1}
\newcommand{\FloatTok}[1]{\textcolor[rgb]{0.25,0.63,0.44}{#1}}
\newcommand{\FunctionTok}[1]{\textcolor[rgb]{0.02,0.16,0.49}{#1}}
\newcommand{\ImportTok}[1]{\textcolor[rgb]{0.00,0.50,0.00}{\textbf{#1}}}
\newcommand{\InformationTok}[1]{\textcolor[rgb]{0.38,0.63,0.69}{\textbf{\textit{#1}}}}
\newcommand{\KeywordTok}[1]{\textcolor[rgb]{0.00,0.44,0.13}{\textbf{#1}}}
\newcommand{\NormalTok}[1]{#1}
\newcommand{\OperatorTok}[1]{\textcolor[rgb]{0.40,0.40,0.40}{#1}}
\newcommand{\OtherTok}[1]{\textcolor[rgb]{0.00,0.44,0.13}{#1}}
\newcommand{\PreprocessorTok}[1]{\textcolor[rgb]{0.74,0.48,0.00}{#1}}
\newcommand{\RegionMarkerTok}[1]{#1}
\newcommand{\SpecialCharTok}[1]{\textcolor[rgb]{0.25,0.44,0.63}{#1}}
\newcommand{\SpecialStringTok}[1]{\textcolor[rgb]{0.73,0.40,0.53}{#1}}
\newcommand{\StringTok}[1]{\textcolor[rgb]{0.25,0.44,0.63}{#1}}
\newcommand{\VariableTok}[1]{\textcolor[rgb]{0.10,0.09,0.49}{#1}}
\newcommand{\VerbatimStringTok}[1]{\textcolor[rgb]{0.25,0.44,0.63}{#1}}
\newcommand{\WarningTok}[1]{\textcolor[rgb]{0.38,0.63,0.69}{\textbf{\textit{#1}}}}
\usepackage{longtable,booktabs,array}
\newcounter{none} % for unnumbered tables
\usepackage{calc} % for calculating minipage widths
% Correct order of tables after \paragraph or \subparagraph
\usepackage{etoolbox}
\makeatletter
\patchcmd\longtable{\par}{\if@noskipsec\mbox{}\fi\par}{}{}
\makeatother
% Allow footnotes in longtable head/foot
\IfFileExists{footnotehyper.sty}{\usepackage{footnotehyper}}{\usepackage{footnote}}
\makesavenoteenv{longtable}
\setlength{\emergencystretch}{3em} % prevent overfull lines
\providecommand{\tightlist}{%
  \setlength{\itemsep}{0pt}\setlength{\parskip}{0pt}}
\usepackage{bookmark}
\IfFileExists{xurl.sty}{\usepackage{xurl}}{} % add URL line breaks if available
\urlstyle{same}
\hypersetup{
  hidelinks,
  pdfcreator={LaTeX via pandoc}}

\author{}
\date{}

\begin{document}

\section{Cilk: An Efficient Multithreaded Runtime
System*}\label{cilk-an-efficient-multithreaded-runtime-system}

ROBERT D. BLUMOFE, CHRISTOPHER F. JOERG, BRADLEY C. KUSZMAUL, CHARLES E.
LEISERSON, KEITH H. RANDALL, AND YULI ZHOU

MIT Laboratory for Computer Science, 545 Technology Square, Cambridge MA
02139

December 19, 1995

\section{Abstract}\label{abstract}

Cilk (pronounced ``silk'') is a C-based runtime system for multithreaded
parallel programming. In this paper, we document the efficiency of the
Cilk work-stealing scheduler, both empirically and analytically. We show
that on real and synthetic applications, the ``work'' and
``critical-path length'' of a Cilk computation can be used to model
performance accurately. Consequently, a Cilk programmer can focus on
reducing the computation's work and critical-path length, insulated from
load balancing and other runtime scheduling issues. We also prove that
for the class of ``fully strict'' (well-structured) programs, the Cilk
scheduler achieves space, time, and communication bounds all within a
constant factor of optimal.

The Cilk runtime system currently runs on the Connection Machine CM5
MPP, the Intel Paragon MPP, the Sun Sparcstation SMP, and the Cilk-NOW
network of workstations. Applications written in Cilk include protein
folding, graphic rendering, backtrack search, and the ★Socrates chess
program, which won second prize in the 1995 World Computer Chess
Championship.

\section{1 Introduction}\label{introduction}

Multithreading has become an increasingly popular way to implement
dynamic, highly asynchronous, concurrent programs {[}1, 9, 10, 11, 12,
13, 16, 21, 23, 24, 26, 27, 30, 36, 37, 39, 42, 43{]}. A multithreaded
system provides the programmer with a means to create, synchronize, and
schedule threads. Although the schedulers in many of these runtime
systems seem to perform well in

level 0 level 1 level 2 level 3\\
\textit{[Image omitted: images/8d3f58d3d343589238b6424cd9afdc2ac5c8e3a32089ea4d8ef221652445258f.jpg]}

flowchart

\begin{Shaded}
\begin{Highlighting}[]
\NormalTok{graph TD}
\NormalTok{    A[" "] {-}{-}\textgreater{} B[" "]}
\NormalTok{    B {-}{-}\textgreater{} C[" "]}
\NormalTok{    C {-}{-}\textgreater{} D[" "]}
\NormalTok{    D {-}{-}\textgreater{} E[" "]}
\NormalTok{    E {-}{-}\textgreater{} F[" "]}
\NormalTok{    F {-}{-}\textgreater{} G[" "]}
\NormalTok{    G {-}{-}\textgreater{} H[" "]}
\NormalTok{    H {-}{-}\textgreater{} I[" "]}
\NormalTok{    I {-}{-}\textgreater{} J[" "]}
\NormalTok{    J {-}{-}\textgreater{} K[" "]}
\NormalTok{    K {-}{-}\textgreater{} L[" "]}
\NormalTok{    L {-}{-}\textgreater{} M[" "]}
\NormalTok{    M {-}{-}\textgreater{} N[" "]}
\NormalTok{    N {-}{-}\textgreater{} O[" "]}
\NormalTok{    O {-}{-}\textgreater{} P[" "]}
\NormalTok{    P {-}{-}\textgreater{} Q[" "]}
\NormalTok{    Q {-}{-}\textgreater{} R[" "]}
\NormalTok{    R {-}{-}\textgreater{} S[" "]}
\NormalTok{    S {-}{-}\textgreater{} T[" "]}
\NormalTok{    T {-}{-}\textgreater{} U[" "]}
\NormalTok{    U {-}{-}\textgreater{} V[" "]}
\NormalTok{    V {-}{-}\textgreater{} W[" "]}
\NormalTok{    W {-}{-}\textgreater{} X[" "]}
\NormalTok{    X {-}{-}\textgreater{} Y[" "]}
\NormalTok{    Y {-}{-}\textgreater{} Z[" "]}
\NormalTok{    Z {-}{-}\textgreater{} A}
\NormalTok{    B {-}{-}\textgreater{} C}
\NormalTok{    C {-}{-}\textgreater{} D}
\NormalTok{    D {-}{-}\textgreater{} E}
\NormalTok{    E {-}{-}\textgreater{} F}
\NormalTok{    F {-}{-}\textgreater{} G}
\NormalTok{    G {-}{-}\textgreater{} H}
\NormalTok{    H {-}{-}\textgreater{} I}
\NormalTok{    I {-}{-}\textgreater{} J}
\NormalTok{    J {-}{-}\textgreater{} K}
\NormalTok{    K {-}{-}\textgreater{} L}
\NormalTok{    L {-}{-}\textgreater{} M}
\NormalTok{    M {-}{-}\textgreater{} N}
\NormalTok{    N {-}{-}\textgreater{} O}
\NormalTok{    O {-}{-}\textgreater{} P}
\NormalTok{    P {-}{-}\textgreater{} Q}
\NormalTok{    Q {-}{-}\textgreater{} R}
\NormalTok{    R {-}{-}\textgreater{} S}
\NormalTok{    S {-}{-}\textgreater{} T}
\NormalTok{    T {-}{-}\textgreater{} U}
\NormalTok{    U {-}{-}\textgreater{} V}
\NormalTok{    V {-}{-}\textgreater{} W}
\NormalTok{    W {-}{-}\textgreater{} X}
\NormalTok{    X {-}{-}\textgreater{} Y}
\NormalTok{    Y {-}{-}\textgreater{} Z}
\NormalTok{    Z {-}{-}\textgreater{} A}
\NormalTok{    style A fill:\#ccc}
\NormalTok{    style B fill:\#ccc}
\NormalTok{    style C fill:\#ccc}
\NormalTok{    style D fill:\#ccc}
\NormalTok{    style E fill:\#ccc}
\NormalTok{    style F fill:\#ccc}
\NormalTok{    style G fill:\#ccc}
\NormalTok{    style H fill:\#ccc}
\NormalTok{    style I fill:\#ccc}
\NormalTok{    style J fill:\#ccc}
\NormalTok{    style K fill:\#ccc}
\NormalTok{    style L fill:\#ccc}
\NormalTok{    style M fill:\#ccc}
\NormalTok{    style N fill:\#ccc}
\NormalTok{    style O fill:\#ccc}
\NormalTok{    style P fill:\#ccc}
\NormalTok{    style Q fill:\#ccc}
\NormalTok{    style R fill:\#ccc}
\NormalTok{    style S fill:\#ccc}
\NormalTok{    style T fill:\#ccc}
\NormalTok{    style U fill:\#ccc}
\NormalTok{    style V fill:\#ccc}
\NormalTok{    style W fill:\#ccc}
\NormalTok{    style X fill:\#ccc}
\NormalTok{    style Y fill:\#ccc}
\NormalTok{    style Z fill:\#ccc}
\end{Highlighting}
\end{Shaded}

Figure 1: The Cilk model of multithreaded computation. Threads are shown
as circles, which are grouped into procedures. Each downward edge
corresponds to a spawn of a child, each horizontal edge corresponds to a
spawn of a successor, and each upward, curved edge corresponds to a data
dependency. The levels of procedures in the spawn tree are indicated at
the left.

practice, none provide users with a guarantee of application
performance. Cilk is a runtime system whose work-stealing scheduler is
efficient in theory as well as in practice. Moreover, it gives the user
an algorithmic model of application performance based on the measures of
``work'' and ``critical-path length'' which can be used to predict the
runtime of a Cilk program accurately.

A Cilk multithreaded computation can be viewed as a directed acyclic
graph (dag) that unfolds dynamically, as is shown schematically in
Figure 1. A Cilk program consists of a collection of Cilk procedures,
each of which is broken into a sequence of threads, which form the
vertices of the dag. Each thread is a nonblocking C function, which
means that it can run to completion without waiting or suspending once
it has been invoked. As one of the threads from a Cilk procedure runs,
it can spawn a child thread which begins a new child procedure. In the
figure, downward edges connect threads and their procedures with the
children they have spawned. A spawn is like a subroutine call, except
that the calling thread may execute concurrently with its child,
possibly spawning additional children. Since threads cannot block in the
Cilk model, a thread cannot spawn children and then wait for values to
be returned. Rather, the thread must additionally spawn a successor
thread to receive the children's return values when they are produced. A
thread and its successors are considered to be parts of the same Cilk
procedure. In the figure, sequences of successor threads that form Cilk
procedures are connected by horizontal edges. Return values, and other
values sent from one thread to another, induce data dependencies among
the threads, where a thread receiving a value cannot begin until another
thread sends the value. Data dependencies are shown as upward, curved
edges in the figure. Thus, a Cilk computation unfolds as a spawn tree
composed of procedures and the spawn edges that connect them to their
children, but the execution is constrained to follow the precedence
relation determined by the dag of threads.

The execution time of any Cilk program on a parallel computer with P
processors is constrained by two parameters of the computation: the work
and the critical-path length. The work, denoted \(T_{1}\) , is the time
used by a one-processor execution of the program, which corresponds to
the sum of the execution times of all the threads. The critical-path
length, denoted \(T_{\infty}\) , is the total amount of time required by
an ideal infinite-processor execution, which corresponds to the largest
sum of thread execution times along any path. With P processors, the
execution time cannot be less than \(T_{1}/P\) or less than
\(T_{\infty}\) . The Cilk scheduler uses ``work stealing'' {[}4, 8, 14,
15, 16, 21, 29, 30, 31, 37, 43{]} to achieve execution time very near to
the sum of these two measures. Off-line techniques for computing such
efficient schedules have been known for a long time {[}6, 18, 19{]}, but
this efficiency

has been difficult to achieve on-line in a distributed environment while
simultaneously using small amounts of space and communication.

We demonstrate the efficiency of the Cilk scheduler both empirically and
analytically. Empirically, we have been able to document that Cilk works
well for dynamic, asynchronous, tree-like, MIMD-style computations. To
date, the applications we have programmed include protein folding,
graphic rendering, backtrack search, and the \(\star\) Socrates chess
program, which won second prize in the 1995 World Computer Chess
Championship running on Sandia National Laboratories' 1824-node Intel
Paragon MPP. Many of these applications pose problems for more
traditional parallel environments, such as message passing {[}41{]} and
data parallel {[}2, 22{]}, because of the unpredictability of the
dynamic workloads on processors. Analytically, we prove that for ``fully
strict'' (well-structured) programs, Cilk's work-stealing scheduler
achieves execution space, time, and communication bounds all within a
constant factor of optimal. To date, all of the applications that we
have coded are fully strict.

The Cilk language extends C with primitives to express parallelism, and
the Cilk runtime system maps the expressed parallelism into a parallel
execution. A Cilk program is preprocessed to C using our cilk2c
type-checking preprocessor {[}34{]}. The C code is then compiled and
linked with a runtime library to run on the target platform. Currently
supported targets include the Connection Machine CM5 MPP, the Intel
Paragon MPP, the Sun Sparcstation SMP, and the Cilk-NOW {[}3, 5{]}
network of workstations. In this paper, we focus on the Connection
Machine CM5 implementation of Cilk. The Cilk scheduler on the CM5 is
written in about 40 pages of C, and it performs communication among
processors using the Strata {[}7{]} active-message library. This
description of the Cilk environment corresponds to the version as of
March 1995.

The remainder of this paper is organized as follows. Section 2 describes
Cilk's runtime data structures and the C language extensions that are
used for programming. Section 3 describes the work-stealing scheduler.
Section 4 documents the performance of several Cilk applications.
Section 5 shows how the work and critical-path length of a Cilk
computation can be used to model performance. Section 6 shows
analytically that the scheduler works well. Finally, Section 7 offers
some concluding remarks and describes our plans for the future.

\section{2 The Cilk programming environment and
implementation}\label{the-cilk-programming-environment-and-implementation}

In this section we describe a C language extension that we have
developed to ease the task of coding Cilk programs. We also explain the
basic runtime data structures that Cilk uses.

A Cilk procedure is broken into a sequence of threads. The procedure
itself is not explicitly specified in the language. Rather, it is
defined implicitly in terms of its constituent threads. A thread T is
defined in a manner similar to a C function definition:

thread T (arg-decls \ldots) \{ stmts \ldots{} \}

The cilk2c type-checking preprocessor {[}34{]} translates T into a C
function of one argument and void return type. The one argument is a
pointer to a closure data structure, illustrated in Figure 2, which
holds the arguments for T. A closure consists of a pointer to the C
function for T, a slot for each of the specified arguments, and a join
counter indicating the number of missing arguments that need to be
supplied before T is ready to run. A closure is ready if it has obtained
all of its arguments, and it is waiting if some arguments are missing.
To run a ready closure, the Cilk

\textit{[Image omitted: images/08d59fe2c0159c63823e53bfdceb1619e823fd6d9dc65e0aca2c64c941e09076.jpg]}

flowchart

\begin{Shaded}
\begin{Highlighting}[]
\NormalTok{graph TD}
\NormalTok{    A["waiting closure"] {-}{-}\textgreater{}|x: T1| B["code"]}
\NormalTok{    A {-}{-}\textgreater{}|x: 1| C["ready closure"]}
\NormalTok{    A {-}{-}\textgreater{}|x: 17| C}
\NormalTok{    A {-}{-}\textgreater{}|x: 6| C}
\NormalTok{    B {-}{-}\textgreater{}|arguments| C}
\NormalTok{    C {-}{-}\textgreater{}|y: T2| D["code"]}
\NormalTok{    C {-}{-}\textgreater{}|y: 0| D}
\NormalTok{    C {-}{-}\textgreater{}|x: 1| D}
\NormalTok{    D {-}{-}\textgreater{}|join counters| A}
\end{Highlighting}
\end{Shaded}

Figure 2: The closure data structure.

scheduler invokes the thread as a C function using the closure itself as
its sole argument. Within the code for the thread, the arguments are
copied out of the closure data structure into local variables. The
closure is allocated from a simple runtime heap when it is created, and
it is returned to the heap when the thread terminates.

The Cilk language supports a data type called a continuation, which is
specified by the type modifier keyword cont. A continuation is
essentially a global reference to an empty argument slot of a closure,
implemented as a compound data structure containing a pointer to a
closure and an offset that designates one of the closure's argument
slots. Continuations can be created and passed among threads, which
enables threads to communicate and synchronize with each other.
Continuations are typed with the C data type of the slot in the closure.

At runtime, a thread can spawn a child thread, and hence a child
procedure, by creating a closure for the child. Spawning is specified in
the Cilk language as follows:

spawn T (args \ldots)

This statement creates a child closure, fills in all available
arguments, and initializes the join counter to the number of missing
arguments. Available arguments are specified with ordinary C syntax. To
specify a missing argument, the user specifies a continuation variable
(type cont) preceded by a question mark. For example, if the second
argument is ?k, then Cilk sets the variable k to a continuation that
refers to the second argument slot of the created closure. If the
closure is ready, that is, it has no missing arguments, then spawn
causes the closure to be immediately posted to the scheduler for
execution. In typical applications, child closures are created with no
missing arguments.

To string a sequence of threads together to make a Cilk procedure, each
thread in the procedure is responsible for creating its successor. To
create a successor thread, a thread executes the following statement:

spawn\_next T (args \ldots)

This statement is semantically identical to spawn, but it informs the
scheduler that the new closure should be treated as a successor, as
opposed to a child. Successor closures are usually created with some
missing arguments, which are filled in by values produced by the
children.

\begin{Shaded}
\begin{Highlighting}[]
\NormalTok{thread fib }\OperatorTok{(}\NormalTok{cont }\DataTypeTok{int}\NormalTok{ k}\OperatorTok{,} \DataTypeTok{int}\NormalTok{ n}\OperatorTok{)}
\OperatorTok{\{}    \ControlFlowTok{if} \OperatorTok{(}\NormalTok{n}\OperatorTok{\textless{}}\DecValTok{2}\OperatorTok{)}
\NormalTok{    send\_argument }\OperatorTok{(}\NormalTok{k}\OperatorTok{,}\NormalTok{n}\OperatorTok{)}
    \ControlFlowTok{else}
    \OperatorTok{\{}\NormalTok{    cont }\DataTypeTok{int}\NormalTok{ x}\OperatorTok{,}\NormalTok{ y}\OperatorTok{;}
\NormalTok{    spawn\_next sum }\OperatorTok{(}\NormalTok{k}\OperatorTok{,} \OperatorTok{?}\NormalTok{x}\OperatorTok{,} \OperatorTok{?}\NormalTok{y}\OperatorTok{);}
\NormalTok{    spawn fib }\OperatorTok{(}\NormalTok{x}\OperatorTok{,}\NormalTok{ n}\OperatorTok{{-}}\DecValTok{1}\OperatorTok{);}
\NormalTok{    spawn fib }\OperatorTok{(}\NormalTok{y}\OperatorTok{,}\NormalTok{ n}\OperatorTok{{-}}\DecValTok{2}\OperatorTok{);}
    \OperatorTok{\}}
\OperatorTok{\}}

\NormalTok{thread sum }\OperatorTok{(}\NormalTok{cont }\DataTypeTok{int}\NormalTok{ k}\OperatorTok{,} \DataTypeTok{int}\NormalTok{ x}\OperatorTok{,} \DataTypeTok{int}\NormalTok{ y}\OperatorTok{)}
\OperatorTok{\{}\NormalTok{    send\_argument }\OperatorTok{(}\NormalTok{k}\OperatorTok{,}\NormalTok{ x}\OperatorTok{+}\NormalTok{y}\OperatorTok{);}
\OperatorTok{\}} 
\end{Highlighting}
\end{Shaded}

Figure 3: A Cilk procedure, consisting of two threads, to compute the
nth Fibonacci number.

A Cilk procedure does not ever return values in the normal way to a
parent procedure. Instead, the programmer must code the parent procedure
as two threads. The first thread spawns the child procedure, passing it
a continuation pointing to the successor thread's closure. The child
sends its ``return'' value explicitly as an argument to the waiting
successor. This strategy of communicating between threads is called
explicit continuation passing. Cilk provides primitives of the following
form to send values from one closure to another:

\begin{Shaded}
\begin{Highlighting}[]
\NormalTok{send\_argument (k, value) }
\end{Highlighting}
\end{Shaded}

This statement sends the value value to the argument slot of a waiting
closure specified by the continuation k. The types of the continuation
and the value must be compatible. The join counter of the waiting
closure is decremented, and if it becomes zero, then the closure is
ready and is posted to the scheduler.

Figure 3 shows the familiar recursive Fibonacci procedure written in
Cilk. It consists of two threads, fib and its successor sum. Reflecting
the explicit continuation-passing style that Cilk supports, the first
argument to each thread is the continuation specifying where the
``return'' value should be placed.

When the fib function is invoked, it first checks to see if the boundary
case has been reached, in which case it uses send\_argument to
``return'' the value of n to the slot specified by continuation k.
Otherwise, it spawns the successor thread sum, as well as two children
to compute the two subcases. Each of these two children is given a
continuation specifying to which argument in the sum thread it should
send its result. The sum thread simply adds the two arguments when they
arrive and sends this result to the slot designated by k.

Although writing in explicit continuation-passing style is somewhat
onerous for the programmer, the decision to break procedures into
separate nonblocking threads simplifies the Cilk runtime system. Each
Cilk thread leaves the C runtime stack empty when it completes. Thus,
Cilk can run on top of a vanilla C runtime system. A common alternative
{[}21, 27, 35, 37{]} is to support a programming style in which a thread
suspends whenever it discovers that required values have not yet been
computed, resuming when the values become available. When a thread
suspends, however, it

may leave temporary values on the runtime stack which must be saved, or
each thread must have its own runtime stack. Consequently, this
alternative strategy requires that the runtime system employs either
multiple stacks or a mechanism to save these temporaries in
heap-allocated storage. Another advantage of Cilk's strategy is that it
allows multiple children to be spawned from a single nonblocking thread,
which saves on context switching. In Cilk, r children can be spawned and
executed with only \(r + 1\) context switches, whereas the alternative
of suspending whenever a thread is spawned causes 2r context switches.
Since our primary interest is in understanding how to build efficient
multithreaded runtime systems, but without redesigning the basic C
runtime system, we chose the alternative of burdening the programmer
with a requirement which is perhaps less elegant linguistically, but
which yields a simple and portable runtime implementation.

Cilk supports a variety of features that give the programmer greater
control over runtime performance. For example, when the last action of a
thread is to spawn a ready thread, the programmer can use the keyword
tail\_call instead of spawn that produces a ``tail call'' to run the new
thread immediately without invoking the scheduler. Cilk also allows
arrays and subarrays to be passed as arguments to closures. Other
features include various abilities to override the scheduler's
decisions, including on which processor a thread should be placed and
how to pack and unpack data when a closure is migrated from one
processor to another.

\section{3 The Cilk work-stealing
scheduler}\label{the-cilk-work-stealing-scheduler}

Cilk's scheduler uses the technique of work stealing {[}4, 8, 14, 15,
16, 21, 29, 30, 31, 37, 43{]} in which a processor (the thief) who runs
out of work selects another processor (the victim) from whom to steal
work, and then steals the shallowest ready thread in the victim's spawn
tree. Cilk's strategy is for thieves to choose victims at random {[}4,
29, 40{]}. We shall now present the implementation of Cilk's
work-stealing scheduler.

At runtime, each processor maintains a local ready pool to hold ready
closures. Each closure has an associated level, which corresponds to the
thread's depth in the spawn tree. The closures for the threads in the
root procedure have level 0, the closures for the threads in the root's
child procedures have level 1, and so on. The ready pool is an array,
illustrated in Figure 4, in which the Lth element contains a linked list
of all ready closures having level L.

Cilk begins executing the user program by initializing all ready pools
to be empty, placing the initial root thread into the level-0 list of
Processor 0's pool, and then starting a scheduling loop on each
processor.

At each iteration through the scheduling loop, a processor first checks
to see whether its ready pool is empty. If it is, the processor
commences work stealing, which will be described shortly. Otherwise, the
processor performs the following steps:

\begin{enumerate}
\def\labelenumi{\arabic{enumi}.}
\tightlist
\item
  Remove the closure at the head of the list of the deepest nonempty
  level in the ready pool.\\
\item
  Extract the thread from the closure, and invoke it.
\end{enumerate}

As a thread executes, it may spawn or send arguments to other threads.
When the thread dies, control returns to the scheduling loop which
advances to the next iteration.

When a thread at level L performs a spawn of a child thread T, the
processor executes the following operations:

\begin{enumerate}
\def\labelenumi{\arabic{enumi}.}
\tightlist
\item
  Allocate and initialize a closure for T.
\end{enumerate}

\textit{[Image omitted: images/2d5fd063c6f14121f0be75be8e3dd6789e96269ce0ac3a719ce8a756aaf65ff1.jpg]}

flowchart

Game execution flowchart showing sequential levels with steal and
execution steps, including `next closure to steal' and `next closure to
execute' annotations.

Figure 4: A processor's ready pool. At each iteration through the
scheduling loop, the processor executes the closure at the head of the
deepest nonempty level in the ready pool. If the ready pool is empty,
the processor becomes a thief and steals the closure at the head of the
shallowest nonempty level in the ready pool of a victim processor chosen
uniformly at random.

\begin{enumerate}
\def\labelenumi{\arabic{enumi}.}
\setcounter{enumi}{1}
\tightlist
\item
  Copy the available arguments into the closure, initialize any
  continuations to point to missing arguments, and initialize the join
  counter to the number of missing arguments.\\
\item
  Label the closure with level \(L + 1\) .\\
\item
  If there are no missing arguments, post the closure to the ready pool
  by inserting it at the head of the level- \((L+1)\) list.
\end{enumerate}

Execution of spawn\_next is similar, except that the closure is labeled
with level L and, if it is ready, posted to the level-L list.

When a thread performs a send\_argument (k, value), the processor
executes the following operations:

\begin{enumerate}
\def\labelenumi{\arabic{enumi}.}
\tightlist
\item
  Find the closure and argument slot referenced by the continuation k.\\
\item
  Place value in the argument slot, and decrement the join counter of
  the closure.\\
\item
  If the join counter goes to zero, post the closure to the ready pool
  at the appropriate level.
\end{enumerate}

When the continuation k refers to a closure on a remote processor,
network communication ensues. The processor that initiated the
send\_argument function sends a message to the remote processor to
perform the operations. The only subtlety occurs in Step 3. If the
closure must be posted, it is posted to the ready pool of the initiating
processor, rather than to that of the remote processor. This policy is
necessary for the scheduler to be provably efficient, but as a practical
matter, we have also had success with posting the closure to the remote
processor's pool.

If a processor begins an iteration of the scheduling loop and finds that
its ready pool is empty, the processor becomes a thief and commences
work stealing as follows:

\begin{enumerate}
\def\labelenumi{\arabic{enumi}.}
\tightlist
\item
  Select a victim processor uniformly at random.\\
\item
  If the victim's ready pool is empty, go back to Step 1.\\
\item
  If the victim's ready pool is nonempty, extract the closure from the
  head of the list in the shallowest nonempty level of the ready pool,
  and execute it.
\end{enumerate}

Work stealing is implemented with a simple request-reply communication
protocol between the thief and victim.

Why steal work from the shallowest level of the ready pool? The reason
is two-fold---one heuristic and one algorithmic. First, to lower
communication costs, we would like to steal large amounts of work, and
in a tree-structured computation, shallow threads are likely to spawn
more work than deep ones. This heuristic notion is the justification
cited by earlier researchers {[}8, 15, 21, 35, 43{]} who proposed
stealing work that is shallow in the spawn tree. We cannot, however,
prove that shallow threads are more likely to spawn work than deep ones.
What we prove in Section 6 is the following algorithmic property. The
threads that are on the ``critical path'' in the dag, are always at the
shallowest level of a processor's ready pool. Consequently, if
processors are idle, the work they steal makes progress along the
critical path.

\section{4 Performance of Cilk
applications}\label{performance-of-cilk-applications}

The Cilk runtime system executes Cilk applications efficiently and with
predictable performance. Specifically, for dynamic, asynchronous,
tree-like applications, Cilk's work-stealing scheduler produces near
optimal parallel speedup while using small amounts of space and
communication. Furthermore, Cilk application performance can be modeled
accurately as a simple function of work

and critical-path length. In this section, we empirically demonstrate
these facts by experimenting with several applications. This section
begins with a look at these applications and then proceeds with a look
at the performance of these applications. In the next section we look at
application performance modeling. The empirical results of this section
and the next confirm the analytical results of Section 6.

\section{Cilk applications}\label{cilk-applications}

We experimented with the Cilk runtime system using several applications,
some synthetic and some real. The applications are described below:

\begin{itemize}
\tightlist
\item
  fib(n) is the same as was presented in Section 2, except that the
  second recursive spawn is replaced by a tail\_call that avoids the
  scheduler. This program is a good measure of Cilk overhead, because
  the thread length is so small.\\
\item
  queens (n) is a backtrack search program that solves the problem of
  placing n queens on a \(n \times n\) chessboard so that no two queens
  attack each other. The Cilk program is based on serial code by R.
  Sargent of the MIT Media Laboratory. Thread length was enhanced by
  serializing the bottom 7 levels of the search tree.\\
\item
  pfold(x, y, z) is a protein-folding program that finds hamiltonian
  paths in a three-dimensional grid of size \(x \times y \times z\)
  using backtrack search {[}38{]}. Written by Chris Joerg of MIT's
  Laboratory for Computer Science and V. Pande of MIT's Center for
  Material Sciences and Engineering, pfold was the first program to
  enumerate all hamiltonian paths in a \(3 \times 4 \times 4\) grid. For
  our experiments, we timed the enumeration of all paths starting with a
  certain sequence.\\
\item
  ray(x, y) is a parallel program for graphics rendering based on the
  serial POV-Ray program, which uses a ray-tracing algorithm. The entire
  POV-Ray system contains over 20,000 lines of C code, but the core of
  POV-Ray is a simple doubly nested loop that iterates over each pixel
  in a two-dimensional image of size \(x \times y\) . For ray we
  converted the nested loops into a 4-ary divide-and-conquer control
  structure using spawns. \(^{1}\) Our measurements do not include the
  approximately 2.4 seconds of startup time required to read and process
  the scene description file.\\
\item
  knary(n, k, r) is a synthetic benchmark whose parameters can be set to
  produce a variety of values for work and critical-path length. It
  generates a tree of depth n and branching factor k in which the first
  r children at every level are executed serially and the remainder are
  executed in parallel. At each node of the tree, the program runs an
  empty ``for'' loop for 400 iterations.\\
\item
  \(\star\) Socrates is a parallel chess program that uses the Jamboree
  search algorithm {[}25, 31{]} to parallelize a minmax tree search. The
  work of the algorithm varies with the number of processors, because it
  does speculative work that may be aborted during runtime. \(\star\)
  Socrates won second prize in the 1995 ICCA World Computer Chess
  Championship running on the 1824-node Intel Paragon at Sandia National
  Laboratories.
\end{itemize}

\textit{[Image omitted: images/25be69216af759dab000253a17d52cbe0da8df37f08a377aeb8433842969ee0c.jpg]}

natural\_image

Black-and-white artistic rendering of a fish swimming in water, creating
ripples (no text or symbols)

\begin{enumerate}
\def\labelenumi{(\alph{enumi})}
\tightlist
\item
  Ray-traced image.
\end{enumerate}

\textit{[Image omitted: images/168697f758c8954765db19ee2585faebd266ad9220fa6020bfbc607a25691d19.jpg]}

natural\_image

Abstract pink and white graphic with a central circular shape and
concentric rings, set against a dark background with scattered dots (no
text or symbols)

\begin{enumerate}
\def\labelenumi{(\alph{enumi})}
\setcounter{enumi}{1}
\tightlist
\item
  Work at each pixel.\\
  Figure 5: (a) An image rendered with the ray program. (b) This image
  shows the amount of time ray took to compute each pixel value. The
  whiter the pixel, the longer ray worked to compute the corresponding
  pixel value.
\end{enumerate}

Many of these applications place heavy demands on the runtime system due
to their dynamic and irregular nature. For example, in the case of
queens and pfold, the size and shape of the backtrack-search tree cannot
be determined without actually performing the search, and the shape of
the tree often turns out to be highly irregular. With speculative work
that may be aborted, the \(\star\) Socrates minmax tree carries this
dynamic and irregular structure to the extreme. In the case of ray, the
amount of time it takes to compute the color of a pixel in an image is
hard to predict and may vary widely from pixel to pixel, as illustrated
in Figure 5. In all of these cases, high performance demands efficient,
dynamic load balancing at runtime.

All experiments were run on a CM5 supercomputer. The CM5 is a massively
parallel computer based on 32MHz SPARC processors with a fat-tree
interconnection network {[}32{]}. The Cilk runtime system on the CM5
performs communication among processors using the Strata {[}7{]}
active-message library.

\section{Application performance}\label{application-performance}

By running our applications and measuring a suite of performance
parameters, we empirically answer many questions about the effectiveness
of the Cilk runtime system. We focus on the following questions. How
efficiently does the runtime system implement the language primitives?
As we add processors, how much faster will the program run? How much
more space will it require? How much more communication will it perform?
We show that for dynamic, asynchronous, tree-like programs, the Cilk
runtime system efficiently implements the language primitives, and that
it is simultaneously efficient with respect to time, space, and
communication. In Section 6, we reach the same conclusion by analytic
means, but in this section we focus on empirical data from the execution
of our Cilk programs.

The execution of a Cilk program with a given set of inputs grows a Cilk
computation that consists of a tree of procedures and a dag of threads.
These structures were introduced in Section 1. We benchmark our
applications with respect to work and critical-path length.

Recall that the work, denoted by \(T_{1}\) , is the time to execute the
Cilk computation on one processor, which corresponds to the sum of the
execution times of all the threads in the dag. The method used to
measure \(T_{1}\) depends on whether the program is deterministic. For
deterministic programs, the computation only depends on the program and
its inputs, and hence, it is independent of the number of processors and
runtime scheduling decisions. \(^{2}\) All of our applications, except
\(\star\) Socrates, are deterministic. For these deterministic
applications, the work performed by any \(P\) -processor run of the
program is equal to the work performed by a 1-processor run (with the
same input values), so we measure the work \(T_{1}\) directly by timing
the 1-processor run. The \(\star\) Socrates program, on the other hand,
uses speculative execution, and therefore, the computation depends on
the number of processors and scheduling decisions made at runtime. In
this case, timing a 1-processor run is not a reasonable way to measure
the work performed by a run with more processors. We must realize that
the work \(T_{1}\) of an execution with \(P\) processors is defined as
the time it takes 1-processor to execute the same computation, not the
same program (with the same inputs). For \(\star\) Socrates we estimate
the work of a \(P\) -processor run by performing the \(P\) -processor
run and timing the execution of every thread and summing. This method
yields an underestimate, since it does not include scheduling costs. In
either case, a \(P\) -processor execution of a Cilk computation with
work \(T_{1}\) must take time at least \(T_{1}/P\) . \(^{3}\) A \(P\)
-processor execution that takes time equal to this \(T_{1}/P\) lower
bound is said to achieve perfect linear speedup.

Recall that the critical-path length, denoted by \(T_{\infty}\) , is the
time to execute the Cilk computation with infinitely many processors,
which corresponds to the largest sum of thread execution times along any
path in the dag. Cilk can measure critical-path length by timestamping
each thread in the dag with the earliest time at which it could have
been executed. Specifically this timestamp is the maximum of the
earliest time that the thread could have been spawned and, for each
argument, the earliest time that the argument could have been sent.
These values, in turn, are computed from the timestamp of the thread
that performed the spawn or sent the argument. In particular, if a
thread performs a spawn, then the earliest time that the spawn could
occur is equal to the earliest time at which the thread could have been
executed (its timestamp) plus the amount of time the thread ran for
until it performed the spawn. The same property holds for the earliest
time that an argument could be sent. The initial thread of the
computation is timestamped zero, and the critical-path length is then
computed as the maximum over all threads of its timestamp plus the
amount of time it executes for. The measured critical-path length does
not include scheduling and communication costs. A \(P\) -processor
execution of a Cilk computation must take at least as long as the
computation's critical-path length \(T_{\infty}\) . Thus, if
\(T_{\infty}\) exceeds \(T_{1}/P\) , then perfect linear speedup cannot
be achieved.

Figure 6 is a table showing typical performance measures for our Cilk
applications. Each column presents data from a single run of a benchmark
application. We adopt the following notations, which are used in the
table. For each application, we have an efficient serial C
implementation, compiled using gcc -02, whose measured runtime is
denoted \(T_{serial}\) . The Cilk computation's

fib(33)

queens(15)

pfold(3,3,4)

ray(500,500)

knary(10,5,2)

knary(10,4,1)

*Socrates(depth 10)(32 proc.)

*Socrates(depth 10)(256 proc)

(computation parameters)

\(T_{serial}\)

8.487

252.1

615.15

729.2

288.6

40.993

1665

1665

\(T_1\)

73.16

254.6

647.8

732.5

314.6

45.43

3644

7023

\(T_{serial}/T_1\)

0.116

0.9902

0.9496

0.9955

0.9174

0.9023

0.4569

0.2371

\(T_\infty\)

0.000326

0.0345

0.04354

0.0415

4.458

0.255

3.134

3.24

\(T_1/T_\infty\)

224417

7380

14879

17650

70.56

178.2

1163

2168

threads

17,108,660

210,740

9,515,098

424,475

5,859,374

873,812

26,151,774

51,685,823

thread length

4.276μs

1208μs

68.08μs

1726μs

53.69μs

51.99μs

139.3μs

135.9μs

(32-processor experiments)

\(T_P\)

2.298

8.012

20.26

21.68

15.13

1.633

126.1

-

\(T_1/P+T_\infty\)

2.287

7.991

20.29

22.93

14.28

1.675

117.0

-

\(T_1/T_P\)

31.84

31.78

31.97

33.79

20.78

27.81

28.90

-

\(T_1/(P\cdot T_P)\)

0.9951

0.9930

0.9992

1.0558

0.6495

0.8692

0.9030

-

space/proc.

70

95

47

39

41

42

386

-

requests/proc.

185.8

48.0

88.6

218.1

92639

3127

23484

-

steals/proc.

56.63

18.47

26.06

79.25

18031

1034

2395

-

(256-processor experiments)

\(T_P\)

0.2892

1.045

2.590

2.765

8.590

0.4636

-

34.32

\(T_1/P+T_\infty\)

0.2861

1.029

2.574

2.903

5.687

0.4325

-

30.67

\(T_1/T_P\)

253.0

243.7

250.1

265.0

36.62

98.00

-

204.6

\(T_1/(P\cdot T_P)\)

0.9882

0.9519

0.9771

1.035

0.1431

0.3828

-

0.7993

space/proc.

66

76

47

32

48

40

-

405

requests/proc.

73.66

80.40

97.79

82.75

151803

7527

-

30646

steals/proc.

24.10

21.20

23.05

18.34

6378

550

-

1540

Figure 6: Performance of Cilk on various applications. All times are in
seconds, except where noted.

work \(T_{1}\) and critical-path length \(T_{\infty}\) are measured on
the CM5 as described above. The measured execution time of the Cilk
program running on P processors of the CM5 is given by \(T_{P}\) . The
row labeled ``threads'' indicates the number of threads executed, and
``thread length'' is the average thread length (work divided by the
number of threads).

Certain derived parameters are also displayed in the table. The ratio
\(T_{serial}/T_{1}\) is the efficiency of the Cilk program relative to
the C program. The ratio \(T_{1}/T_{\infty}\) is the average
parallelism. The value \(T_{1}/P + T_{\infty}\) is a simple model of the
runtime, which will be discussed later. The speedup is \(T_{1}/T_{P}\) ,
and the parallel efficiency is \(T_{1}/(P \cdot T_{P})\) . The row
labeled ``space/proc.'' indicates the maximum number of closures
allocated at any time on any processor. The row labeled
``requests/proc.'' indicates the average number of steal requests made
by a processor during the execution, and ``steals/proc.'' gives the
average number of closures actually stolen.

The data in Figure 6 shows two important relationships: one between
efficiency and thread length, and another between speedup and average
parallelism.

Considering the relationship between efficiency \(T_{serial}/T_{1}\) and
thread length, we see that for programs with moderately long threads,
the Cilk runtime system induces little overhead. The queens, pfold, ray,
and knary programs have threads with average length greater than 50
microseconds and have efficiency greater than 90 percent. On the other
hand, the fib program has low efficiency, because the threads are so
short: fib does almost nothing besides spawn and send\_argument.

Despite it's long threads, the \(\star\) Socrates program has low
efficiency, because its parallel Jamboree search algorithm is based on
speculatively searching subtrees that are not searched by a serial
algorithm. Consequently, as we increase the number of processors, the
program executes more threads and, hence, does more work. For example,
the 256-processor execution did 7023 seconds of work whereas the
32-processor execution did only 3644 seconds of work. Both of these
executions did considerably more work than the serial program's 1665
seconds of work. Thus, although we observe low efficiency, it is due to
the parallel algorithm and not to Cilk overhead.

Looking at the speedup \(T_1 / T_P\) measured on 32 and 256 processors,
we see that when the average parallelism \(T_1 / T_\infty\) is large
compared with the number \(P\) of processors, Cilk programs achieve
nearly perfect linear speedup, but when the average parallelism is
small, the speedup is much less. The fib, queens, pfold, and ray
programs, for example, have in excess of 7000-fold parallelism and
achieve more than 99 percent of perfect linear speedup on 32 processors
and more than 95 percent of perfect linear speedup on 256 processors.
The \(\star\) Socrates program exhibits somewhat less parallelism and
also somewhat less speedup. On 32 processors the \(\star\) Socrates
program has 1163-fold parallelism, yielding 90 percent of perfect linear
speedup, while on 256 processors it has 2168-fold parallelism yielding
80 percent of perfect linear speedup. With even less parallelism, as
exhibited in the knary benchmarks, less speedup is obtained. For
example, the knary (10, 5, 2) benchmark exhibits only 70-fold
parallelism, and it realizes barely more than 20-fold speedup on 32
processors (less than 65 percent of perfect linear speedup). With
178-fold parallelism, knary (10, 4, 1) achieves 27-fold speedup on 32
processors (87 percent of perfect linear speedup), but only 98-fold
speedup on 256 processors (38 percent of perfect linear speedup).

Although these speedup measures reflect the Cilk scheduler's ability to
exploit parallelism, to obtain application speedup, we must factor in
the efficiency of the Cilk program compared

with the serial C program. Specifically, the application speedup
\(T_{serial}/T_{P}\) is the product of efficiency \(T_{serial}/T_{1}\)
and speedup \(T_{1}/T_{P}\) . For example, applications such as fib and
\(\star\) Socrates with low efficiency generate correspondingly low
application speedup. The \(\star\) Socrates program, with efficiency
0.2371 and speedup 204.6 on 256 processors, exhibits application speedup
of \(0.2371 \cdot 204.6 = 48.51\) . For the purpose of understanding
scheduler performance, we have decoupled the efficiency of the
application from the efficiency of the scheduler.

Looking more carefully at the cost of a spawn in Cilk, we find that it
takes a fixed overhead of about 50 cycles to allocate and initialize a
closure, plus about 8 cycles for each word argument. In comparison, a C
function call on a CM5 SPARC processor takes 2 cycles of fixed overhead
(assuming no register window overflow) plus 1 cycle for each word
argument (assuming all arguments are transferred in registers). Thus, a
spawn in Cilk is roughly an order of magnitude more expensive than a C
function call. This Cilk overhead is quite apparent in the fib program,
which does almost nothing besides spawn and send\_argument. Based on
fib's measured efficiency of 0.116, we can conclude that the aggregate
average cost of a spawn/send\_argument in Cilk is between 8 and 9 times
the cost of a function call/return in C.

Efficient execution of programs with short threads requires a
low-overhead spawn operation. As can be observed from Figure 6, the vast
majority of threads execute on the same processor on which they are
spawned. For example, the fib program executed over 17 million threads
but migrated only 6170 (24.10 per processor) when run with 256
processors. Taking advantage of this property, other researchers {[}17,
27, 35{]} have developed techniques for implementing spawns such that
when the child thread executes on the same processor as its parent, the
cost of the spawn operation is roughly equal the cost of a function
call. We hope to incorporate such techniques into future implementations
of Cilk.

Finally, we make two observations about the space and communication
measures in Figure 6.

Looking at the ``space/proc.'' rows, we observe that the space per
processor is generally quite small and does not grow with the number of
processors. For example, \(\star\) Socrates on 32 processors executes
over 26 million threads, yet no processor ever contains more than 386
allocated closures. On 256 processors, the number of executed threads
nearly doubles to over 51 million, but the space per processor barely
changes. In Section 6 we show formally that for an important class of
Cilk programs, the space per processor does not grow as we add
processors.

Looking at the ``requests/proc.'' and ``steals/proc.'' rows in Figure 6,
we observe that the amount of communication grows with the critical-path
length but does not grow with the work. For example, fib, queens, pfold,
and ray all have critical-path lengths under a tenth of a second long
and perform fewer than 220 requests and 80 steals per processor, whereas
knary(10, 5, 2) and \(\star\) Socrates have critical-path lengths more
than 3 seconds long and perform more than 20,000 requests and 1500
steals per processor. The table does not show any clear correlation
between work and either requests or steals. For example, ray does more
than twice as much work as knary(10, 5, 2), yet it performs two orders
of magnitude fewer requests. In Section 6, we show that for a class of
Cilk programs, the communication per processor grows at most linearly
with the critical-path length and does not grow as a function of the
work.

\section{5 Modeling performance}\label{modeling-performance}

We further document the effectiveness of the Cilk scheduler by showing
empirically that Cilk application performance can be modeled accurately
with a simple function of work \(T_{1}\) and critical-path length
\(T_{\infty}\) . Specifically, we use the knary synthetic benchmark to
show that the runtime of an application on P processors can be modeled
as \(T_{P} \approx T_{1}/P + c_{\infty}T_{\infty}\) , where
\(c_{\infty}\) is a small constant (about 1.5 for knary) determined by
curve fitting. This result shows that we obtain nearly perfect linear
speedup when the critical path is short compared with the average amount
of work per processor. We also show that a model of this kind is
accurate even for \(\star\) Socrates, which is our most complex
application programmed to date.

We would like our scheduler to execute a Cilk computation with \(T_{1}\)
work in \(T_{1}/P\) time on \(P\) processors. Such perfect linear
speedup cannot be obtained whenever the computation's critical-path
length \(T_{\infty}\) exceeds \(T_{1}/P\) , since we always have
\(T_{P} \geq T_{\infty}\) or more generally,
\(T_{P} \geq \max \{T_{1}/P, T_{\infty}\}\) . The critical-path length
\(T_{\infty}\) is the stronger lower bound on \(T_{P}\) whenever \(P\)
exceeds the average parallelism \(T_{1}/T_{\infty}\) , and \(T_{1}/P\)
is the stronger bound otherwise. A good scheduler should meet each of
these bounds as closely as possible.

In order to investigate how well the Cilk scheduler meets these two
lower bounds, we used our synthetic knary benchmark, which can grow
computations that exhibit a range of values for work and critical-path
length.

Figure 7 shows the outcome from many experiments of running knary with
various input values (n, k, and r) on various numbers of processors. The
figure plots the measured speedup \(T_{1} / T_{P}\) for each run against
the machine size \(P\) for that run. In order to compare the outcomes
for runs with different input values, we have normalized the plotted
value for each run as follows. In addition to the speedup, we measure
for each run the work \(T_{1}\) and the critical-path length
\(T_{\infty}\) , as previously described. We then normalize the machine
size and the speedup by dividing these values by the average parallelism
\(T_{1} / T_{\infty}\) . For each run, the horizontal position of the
plotted datum is \(P / (T_{1} / T_{\infty})\) , and the vertical
position of the plotted datum is
\((T_{1} / T_{P}) / (T_{1} / T_{\infty}) = T_{\infty} / T_{P}\) .
Consequently, on the horizontal axis, the normalized machine size is 1.0
when the number of processors is equal to the average parallelism. On
the vertical axis, the normalized speedup is 1.0 when the runtime equals
the critical-path length. We can draw the two lower bounds on time as
upper bounds on speedup. The horizontal line at 1.0 is the upper bound
on speedup obtained from the critical-path length,
\(T_{P} \geq T_{\infty}\) , and the 45-degree line is the linear speedup
bound, \(T_{P} \geq T_{1} / P\) . As can be seen from the figure, on the
knary runs for which the average parallelism exceeds the number of
processors (normalized machine size less than 1), the Cilk scheduler
obtains nearly perfect linear speedup. In the region where the number of
processors is large compared to the average parallelism (normalized
machine size greater than 1), the data is more scattered, but the
speedup is always within a factor of 4 of the critical-path length upper
bound.

The theoretical results from Section 6 show that the expected running
time of a Cilk computation on P processors is
\(T_{P} = O(T_{1}/P + T_{\infty})\) . Thus, it makes sense to try to fit
the knary data to a curve of the form
\(T_{P} = c_{1}(T_{1}/P) + c_{\infty}(T_{\infty})\) . A least-squares
fit to the data to minimize the relative error yields
\(c_{1} = 0.9543 \pm 0.1775\) and \(c_{\infty} = 1.54 \pm 0.3888\) with
95 percent confidence. The \(R^{2}\) correlation coefficient of the fit
is 0.989101, and the mean relative error is 13.07 percent. The curve fit
is shown in Figure 7, which also plots the simpler curves
\(T_{P} = T_{1}/P + T_{\infty}\) and
\(T_{P} = T_{1}/P + 2 \cdot T_{\infty}\) for comparison. As can be seen
from the figure, little is lost in the linear speedup range of the curve
by assuming that the coefficient \(c_{1}\) on the \(T_{1}/P\) term
equals 1. Indeed,

\textit{[Image omitted: images/c7090950d7ddd8a090aca99b0218ce7295bb04130efb265ea572c5427a04bfaa.jpg]}

line

{\def\LTcaptype{none} % do not increment counter
\begin{longtable}[]{@{}ll@{}}
\toprule\noalign{}
Normalized Machine Size & Normalized Speedup \\
\midrule\noalign{}
\endhead
\bottomrule\noalign{}
\endlastfoot
0.0001 & 0.0001 \\
0.001 & 0.001 \\
0.01 & 0.01 \\
0.1 & 0.1 \\
1 & 1 \\
10 & 1 \\
\end{longtable}
}

Figure 7: Normalized speedups for the knary synthetic benchmark using
from 1 to 256 processors. The horizontal axis is the number P of
processors and the vertical axis is the speedup \(T_{1}/T_{P}\) , but
each data point has been normalized by dividing by \(T_{1}/T_{\infty}\)
.

\textit{[Image omitted: images/47153418a889bb98d8b82b1edb3f1df194767be053eb657bbd0bb1875c890c3d.jpg]}

scatter

{\def\LTcaptype{none} % do not increment counter
\begin{longtable}[]{@{}ll@{}}
\toprule\noalign{}
Normalized Machine Size & Normalized Speedup \\
\midrule\noalign{}
\endhead
\bottomrule\noalign{}
\endlastfoot
0.01 & 0.01 \\
0.1 & 0.1 \\
1 & 1 \\
\end{longtable}
}

Figure 8: Normalized speedups for the \(\star\) Socrates chess program.

a fit to \(T_{P} = T_{1}/P + c_{\infty}(T_{\infty})\) yields
\(c_{\infty} = 1.509 \pm 0.3727\) with \(R^{2} = 0.983592\) and a mean
relative error of 4.04 percent, which is in some ways better than the
fit that includes a \(c_{1}\) term. (The \(R^{2}\) measure is a little
worse, but the mean relative error is much better.)

It makes sense that the data points become more scattered when P is
close to or exceeds the average parallelism. In this range, the amount
of time spent in work stealing becomes a significant fraction of the
overall execution time. The real measure of the quality of a scheduler
is how much larger than P the average parallelism \(T_{1}/T_{\infty}\)
must be before \(T_{P}\) shows substantial influence from the
critical-path length. One can see from Figure 7 that if the average
parallelism exceeds P by a factor of 10, the critical-path length has
almost no impact on the running time.

To confirm our simple model of the Cilk scheduler's performance on a
real application, we ran \(\star\) Socrates on a variety of chess
positions using various numbers of processors. Figure 8 shows the
results of our study, which confirm the results from the knary synthetic
benchmark. The best fit to
\(T_{P} = c_{1}(T_{1} / P) + c_{\infty}(T_{\infty})\) yields
\(c_{1} = 1.067 \pm 0.0141\) and \(c_{\infty} = 1.042 \pm 0.0467\) with
95 percent confidence. The \(R^{2}\) correlation coefficient of the fit
is 0.9994, and the mean relative error is 4.05 percent.

Indeed, as some of us were developing and tuning heuristics to increase
the performance of \(\star\) Socrates, we used work and critical-path
length as our measures of progress. This methodology let us avoid being
trapped by the following interesting anomaly. We made an ``improvement''
that sped up the program on 32 processors. From our measurements,
however, we discovered that it was faster only because it saved on work
at the expense of a much longer critical path. Using the simple model
\(T_{P} = T_{1} / P + T_{\infty}\) , we concluded that on a
512-processor Connection Machine CM5 MPP at the National Center for
Supercomputer Applications at the University of Illinois,
Urbana-Champaign, which was our platform for our early tournaments, the
``improvement'' would yield a loss of performance, a fact that we later
verified. Measuring work and critical-path length enabled us to use
experiments on a 32-processor machine to improve our program for the
512-processor machine, but without using the 512-processor machine, on
which computer time was scarce.

\section{6 A theoretical analysis of the Cilk
scheduler}\label{a-theoretical-analysis-of-the-cilk-scheduler}

In this section we use algorithmic analysis techniques to prove that for
the class of ``fully strict'' Cilk programs, Cilk's work-stealing
scheduling algorithm is efficient with respect to space, time, and
communication. A fully strict program is one for which each thread sends
arguments only to its parent's successor threads. In the analysis and
bounds of this section, we further assume that each thread spawns at
most one successor thread. Programs such as \(\star\) Socrates violate
this assumption, and at the end of the section, we explain how the
analysis and bounds can be generalized to handle such programs. For
fully strict programs, we prove the following three bounds on space,
time, and communication:

Space The space used by a P-processor execution is bounded by
\(S_{P} \leq S_{1}P\) , where \(S_{1}\) denotes the space used by the
serial execution of the Cilk program. This bound is existentially
optimal to within a constant factor {[}4{]}.

Time With P processors, the expected execution time, including
scheduling overhead, is bounded by \(O(T_{1}/P + T_{\infty})\) . Since
both \(T_{1}/P\) and \(T_{\infty}\) are lower bounds for any P-processor
execution, this bound is within a constant factor of optimal.

Communication The expected number of bytes communicated during a
P-processor execution is \(O(PT_{\infty}S_{\max})\) , where \(S_{max}\)
is the size of the largest closure in the computation. This bound is
existentially optimal to within a constant factor {[}44{]}.

The expected-time bound and the expected-communication bound can be
converted into high-probability bounds at the cost of only a small
additive term in both cases. Full proofs of these bounds, using
generalizations of the techniques developed in {[}4{]}, can be found in
{[}3{]}.

The space bound can be obtained from a ``busy-leaves'' property {[}4{]}
that characterizes the allocated closures at all times during the
execution. In order to state this property simply, we first define some
terms. We say that two or more closures are siblings if they were
spawned by the same parent, or if they are successors (by one or more
spawn\_next's) of closures spawned by the same parent. Sibling closures
can be ordered by age: the first child spawned is older than the second,
and so on. At any given time during the execution, we say that a closure
is a leaf if it has no allocated children, and we say that a leaf
closure is a primary leaf if, in addition, it has no younger siblings
allocated. The busy-leaves property states that every primary-leaf
closure has a processor working on it.

\section{Lemma 1 Cilk's scheduler maintains the busy-leaves
property.}\label{lemma-1-cilks-scheduler-maintains-the-busy-leaves-property.}

Proof: Consider the three possible ways that a primary-leaf closure can
be created. First, when a thread spawns children, the youngest of these
children is a primary leaf. Second, when a thread completes and its
closure is freed, if that closure has an older sibling and that sibling
has no children, then the older-sibling closure becomes a primary leaf.
Finally, when a thread completes and its closure is freed, if that
closure has no allocated siblings, then the youngest closure of its
parent's successor threads is a primary leaf. The induction follows by
observing that in all three of these cases, Cilk's scheduler guarantees
that a processor works on the new primary leaf. In the third case we use
the important fact that a newly activated closure is posted on the
processor that activated it (and not on the processor on which it was
residing).

Theorem 2 For any fully strict Cilk program, if \(S_{1}\) is the space
used to execute the program on 1 processor, then with any number P of
processors, Cilk's work-stealing scheduler uses at most \(S_{1}P\)
space.

Proof: We shall obtain the space bound \(S_{P} \leq S_{1}P\) by
assigning every allocated closure to a primary leaf such that the total
space of all closures assigned to a given primary leaf is at most
\(S_{1}\) . Since Lemma 1 guarantees that all primary leaves are busy,
at most P primary-leaf closures can be allocated, and hence the total
amount of space is at most \(S_{1}P\) .

The assignment of allocated closures to primary leaves is made as
follows. If the closure is a primary leaf, it is assigned to itself.
Otherwise, if the closure has any allocated children, then it is
assigned to the same primary leaf as its youngest child. If the closure
is a leaf but has some younger siblings, then the closure is assigned to
the same primary leaf as its youngest sibling. In this recursive
fashion, we assign every allocated closure to a primary leaf. Now, we
consider the set of closures assigned to a given primary leaf. The total
space of these closures is at most \(S_{1}\) , because this set of
closures is a subset of the closures that are allocated during a
1-processor execution when the processor is executing this primary leaf,
which completes the proof.

We are now ready to analyze execution time. Our strategy is to mimic the
theorems of {[}4{]} for a more restricted model of multithreaded
computation. As in {[}4{]}, the bounds assume a communication model in
which messages are delayed only by contention at destination processors,
but no assumptions are made about the order in which contending messages
are delivered {[}33{]}. For technical reasons in our analysis of
execution time, the critical path is calculated assuming that all
threads spawned by a parent thread are spawned at the end of the parent
thread.

In our analysis of execution time, we use an accounting argument. At
each time step, each of the P processors places a dollar in one of three
buckets according to its actions at that step. If the processor executes
an instruction of a thread at the step, it places its dollar into the
WORK bucket. If the processor initiates a steal attempt, it places its
dollar into the STEAL bucket. Finally, if the processor merely waits for
a steal request that is delayed by contention, then it places its dollar
into the WAIT bucket. We shall derive the running time bound by upper
bounding the dollars in each bucket at the end of the computation,
summing these values, and then dividing by P, the total number of
dollars put into buckets on each step.

Lemma 3 When the execution of a fully strict Cilk computation with work
\(T_{1}\) ends, the WORK bucket contains \(T_{1}\) dollars.

Proof: The computation contains a total of \(T_{1}\) instructions.

Lemma 4 When the execution of a fully strict Cilk computation ends, the
expected number of dollars in the WAIT bucket is less than the number of
dollars in the STEAL bucket.

Proof: Lemma 5 of {[}4{]} shows that if P processors make M random steal
requests during the course of a computation, where requests with the
same destination are serially queued at the destination, then the
expected total delay is less than M.

Lemma 5 When the P-processor execution of a fully strict Cilk
computation with critical-path length \(T_{\infty}\) and for which each
thread has at most one successor ends, the expected number of dollars in
the STEAL bucket is \(O(PT_{\infty})\) .

Proof sketch: The proof follows the delay-sequence argument of {[}4{]},
but with some differences that we shall point out. Full details can be
found in {[}3{]}, which generalizes to the situation in which a thread
can have more than one successor.

At any given time during the execution, we say that a thread is critical
if it has not yet been executed but all of its predecessors in the dag
have been executed. For this argument, the dag must be augmented with
``ghost'' threads and additional edges to represent implicit
dependencies imposed by the Cilk scheduler. We define a delay sequence
to be a pair \((P, s)\) such that P is a path of threads in the
augmented dag and s is a positive integer. We say that a delay sequence
\((P, s)\) occurs in an execution if at least s steal attempts are
initiated while some thread of P is critical.

The next step of the proof is to show that if at least s steal attempts
occur during an execution, where s is sufficiently large, then some
delay sequence \((P, s)\) must occur. That is, there must be some path P
in the dag such that each of the s steal attempts occurs while some
thread of P is critical. We do not give the construction here, but
rather refer the reader to {[}3, 4{]} for directly analogous arguments.

The last step of the proof is to show that a delay sequence with
\(s = \Omega(PT_{\infty})\) is unlikely to occur. The key to this step
is a lemma, which describes the structure of threads the processors'
ready pools. This structural lemma implies that if a thread is critical,
it is the next thread to be stolen from the pool in which it resides.
Intuitively, after P steal attempts, we expect one of these attempts to
have targeted the processor in which the critical thread of interest
resides. In this case, the critical thread will be stolen and executed,
unless, of course, it has already been executed by the local processor.
Thus, after \(PT_{\infty}\) steal attempts, we expect all threads on P
to have been executed. The delay-sequence argument formalizes this
intuition. Thus, the expected number s of dollars in the STEAL bucket is
at most \(O(PT_{\infty})\) .

Theorem 6 Consider any fully strict Cilk computation with work \(T_{1}\)
and critical-path length \(T_{\infty}\) such that every thread spawns at
most one successor. With any number P of processors, Cilk's
work-stealing scheduler runs the computation in expected time
\(O(T_{1}/P + T_{\infty})\) .

Proof: We sum the dollars in the three buckets and divide by P. By Lemma
3, the WORK bucket contains \(T_{1}\) dollars. By Lemma 4, the WAIT
bucket contains at most a constant times the number of dollars in the
STEAL bucket, and Lemma 5 implies that the total number of dollars in
both buckets is \(O(PT_{\infty})\) . Thus, the sum of the dollars is
\(T_{1} + O(PT_{\infty})\) , and the bound on execution time is obtained
by dividing by P.

In fact, it can be shown using the techniques of {[}4{]} that for any
\(\epsilon > 0\) , with probability at least \(1 - \epsilon\) , the
execution time on P processors is
\(O(T_{1}/P + T_{\infty} + \lg P + \lg(1/\epsilon))\) .

Theorem 7 Consider any fully strict Cilk computation with work \(T_{1}\)
and critical-path length \(T_{\infty}\) such that every thread spawns at
most one successor. For any number P of processors, the total number of
bytes communicated by Cilk's work-stealing scheduler has expectation
\(O(PT_{\infty}S_{max})\) , where \(S_{max}\) is the size in bytes of
the largest closure in the computation.

Proof: The proof follows directly from Lemma 5. All communication costs
can be associated with steals or steal requests, and at most
\(O(S_{\max})\) bytes are communicated for each successful steal.

In fact, for any \(\epsilon > 0\) , the probability is at least
\(1 - \epsilon\) that the total communication incurred is
\(O(P(T_{\infty} + \lg(1/\epsilon))S_{\max})\) .

The analysis and bounds we have derived apply to fully strict programs
in the case when each thread spawns at most one successor. Some
programs, such as \(\star\) Socrates, contain threads that spawn several
successors. In {[}3{]}, the theorems above are generalized to handle
this situation as follows. Let \(n_l\) denote the maximum number of
threads belonging to any one procedure such that all the threads are
simultaneously living during some execution. Let \(n_d\) denote the
maximum number of dependency edges between any pair of threads. When
each thread can spawn at most one successor, we have \(n_l = 1\) and
\(n_d = 1\) and the theorems as proved in this paper hold. When \(n_l\)
or \(n_d\) exceeds 1, however, the arguments must be modified.

Specifically, when \(n_{l}\) or \(n_{d}\) exceeds 1, the analysis of the
number of dollars in the STEAL bucket must be modified. A critical
thread may no longer be the first thread to be stolen from a processor's
ready pool. Other noncritical threads from the same procedure may be
stolen in advance of the critical thread. Moreover, extra dependency
edges may cause even more steal

attempts to occur before a critical thread gets stolen. Accounting for
these extra steals in the argument, we obtain the following bounds on
time and communication. For any number P of processors, the expected
execution time is \(O(T_{1}/P + n_{l}T_{\infty})\) , and the expected
number of bytes communicated is
\(O(n_{l}PT_{\infty}(n_{d} + S_{\max}))\) . The \(O(S_{1}P)\) bound on
space is unchanged. Analogous high-probability bounds for time and
communication can be found in {[}3{]}.

\section{7 Conclusion}\label{conclusion}

To produce high-performance parallel applications, programmers often
focus on communication costs and execution time, quantities that are
dependent on specific machine configurations. We argue that a programmer
should think instead about work and critical-path length, abstractions
that can be used to characterize the performance of an algorithm
independent of the machine configuration. Cilk provides a programming
model in which work and critical-path length are observable quantities,
and it delivers guaranteed performance as a function of these
quantities. Work and critical-path length have been used in the theory
community for years to analyze parallel algorithms {[}28{]}. Blelloch
{[}2{]} has developed a performance model for data-parallel computations
based on these same two abstract measures. He cites many advantages to
such a model over machine-based models. Cilk provides a similar
performance model for the domain of asynchronous, multi-threaded
computation.

Although Cilk offers performance guarantees, its current capabilities
are limited, and programmers find its explicit continuation-passing
style to be onerous. Cilk is good at expressing and executing dynamic,
asynchronous, tree-like, MIMD computations, but it is not yet ideal for
more traditional parallel applications that can be programmed
effectively in, for example, a message-passing, data-parallel, or
single-thread-per-processor, shared-memory style. We are currently
working on extending Cilk's capabilities to broaden its applicability. A
major constraint is that we do not want new features to destroy Cilk's
guarantees of performance. Our current research focuses on implementing
``dag-consistent'' shared memory, which allows programs to operate on
shared memory without costly communication or hardware support; on
providing a linguistic interface that produces continuation-passing code
for our runtime system from a more traditional call-return specification
of spawns; and on incorporating persistent threads and less strict
semantics in ways that do not destroy the guaranteed performance of our
scheduler. Recent information about Cilk is maintained on the World Wide
Web in page http://theory.lcs.mit.edu/\textasciitilde cilk.

\section{Acknowledgments}\label{acknowledgments}

We gratefully acknowledge the inspiration of Michael Halbherr, now of
the Boston Consulting Group in Zurich, Switzerland. Mike's PCM runtime
system {[}20{]} developed at MIT was the precursor of Cilk, and many of
the design decisions in Cilk are owed to him. We thank Shail Aditya and
Sivan Toledo of MIT and Larry Rudolph of Hebrew University for helpful
discussions. Xinmin Tian of McGill University provided helpful
suggestions for improving the paper. Don Dailey and International Master
Larry Kaufman, both formerly of Heuristic Software, were part of the
\(\star\) Socrates development team. Rolf Riesen of Sandia National
Laboratories ported Cilk to the Intel Paragon MPP running under the
SUNMOS operating system, John Litvin and Mike Stupak ported

Cilk to the Paragon running under OSF, and Andy Shaw of MIT ported Cilk
to SMP platforms. Thanks to Matteo Frigo and Rob Miller of MIT for their
many contributions to the Cilk system. Thanks to the Scout project at
MIT and the National Center for Supercomputing Applications at
University of Illinois, Urbana-Champaign for access to their CM5
supercomputers for running our experiments. Finally, we acknowledge the
influence of Arvind and his dataflow research group at MIT. Their
pioneering work attracted us to this path, and their vision continues to
challenge us.

\section{References}\label{references}

{[}1{]} Anderson, T. E., Bershad, B. N., Lazowska, E. D., and Levy, H.
M. Scheduler activations: Effective kernel support for the user-level
management of parallelism. In Proceedings of the Thirteenth ACM
Symposium on Operating Systems Principles, pp.~95--109, Pacific Grove,
California, Oct.~1991.\\
{[}2{]} Blelloch, G. E. Programming parallel algorithms. In Proceedings
of the 1992 Dartmouth Institute for Advanced Graduate Studies (DAGS)
Symposium on Parallel Computation, pp.~11--18, Hanover, New Hampshire,
Jun.~1992.\\
{[}3{]} Blumofe, R. D. Executing Multithreaded Programs Efficiently.
Ph.D.~thesis, Department of Electrical Engineering and Computer Science,
Massachusetts Institute of Technology, Sep.~1995.\\
{[}4{]} Blumofe, R. D. and Leiserson, C. E. Scheduling multithreaded
computations by work stealing. In Proceedings of the 35th Annual
Symposium on Foundations of Computer Science, pp.~356--368, Santa Fe,
New Mexico, Nov.~1994.\\
{[}5{]} Blumofe, R. D. and Park, D. S. Scheduling large-scale parallel
computations on networks of workstations. In Proceedings of the Third
International Symposium on High Performance Distributed Computing,
pp.~96--105, San Francisco, California, Aug.~1994.\\
{[}6{]} Brent, R. P. The parallel evaluation of general arithmetic
expressions. Journal of the ACM, 21(2):201--206, Apr.~1974.\\
{[}7{]} Brewer, E. A. and Blumofe, R. Strata: A multi-layer
communications library. Technical Report to appear, MIT Laboratory for
Computer Science. Available as
ftp://ftp.lcs.mit.edu/pub/supertech/strata/strata.tar.Z.\\
{[}8{]} Burton, F. W. and Sleep, M. R. Executing functional programs on
a virtual tree of processors. In Proceedings of the 1981 Conference on
Functional Programming Languages and Computer Architecture,
pp.~187--194, Portsmouth, New Hampshire, Oct.~1981.\\
{[}9{]} Carlisle, M. C., Rogers, A., Reppy, J. H., and Hendren, L. J.
Early experiences with Olden. In Proceedings of the Sixth Annual
Workshop on Languages and Compilers for Parallel Computing, Portland,
Oregon, Aug.~1993.\\
{[}10{]} Chandra, R., Gupta, A., and Hennessy, J. L. COOL: An
object-based language for parallel programming. IEEE Computer,
27(8):13--26, Aug.~1994.\\
{[}11{]} Chase, J. S., Amador, F. G., Lazowska, E. D., Levy, H. M., and
Littlefield, R. J. The Amber system: Parallel programming on a network
of multiprocessors. In Proceedings of the Twelfth ACM Symposium on
Operating Systems Principles, pp.~147--158, Litchfield Park, Arizona,
Dec.~1989.

{[}12{]} Cooper, E. C. and Draves, R. P. C Threads. Tech.
Rep.~CMU-CS-88-154, School of Computer Science, Carnegie-Mellon
University, Jun.~1988.\\
{[}13{]} Culler, D. E., Sah, A., Schauser, K. E., von Eicken, T., and
Wawrzynek, J. Fine-grain parallelism with minimal hardware support: A
compiler-controlled threaded abstract machine. In Proceedings of the
Fourth International Conference on Architectural Support for Programming
Languages and Operating Systems, pp.~164--175, Santa Clara, California,
Apr.~1991.\\
{[}14{]} Feldmann, R., Mysliwietz, P., and Monien, B. Studying overheads
in massively parallel min/max-tree evaluation. In Proceedings of the
Sixth Annual ACM Symposium on Parallel Algorithms and Architectures,
pp.~94--103, Cape May, New Jersey, Jun.~1994.\\
{[}15{]} Finkel, R. and Manber, U. DIB---a distributed implementation of
backtracking. ACM Transactions on Programming Languages and Systems,
9(2):235--256, Apr.~1987.\\
{[}16{]} Freeh, V. W., Lowenthal, D. K., and Andrews, G. R. Distributed
Filaments: Efficient fine-grain parallelism on a cluster of
workstations. In Proceedings of the First Symposium on Operating Systems
Design and Implementation, pp.~201--213, Monterey, California,
Nov.~1994.\\
{[}17{]} Goldstein, S. C., Schauser, K. E., and Culler, D. Enabling
primitives for compiling parallel languages. In Third Workshop on
Languages, Compilers, and Run-Time Systems for Scalable Computers, Troy,
New York, May 1995.\\
{[}18{]} Graham, R. L. Bounds for certain multiprocessing anomalies. The
Bell System Technical Journal, 45:1563--1581, Nov.~1966.\\
{[}19{]} Graham, R. L. Bounds on multiprocessing timing anomalies. SIAM
Journal on Applied Mathematics, 17(2):416--429, Mar.~1969.\\
{[}20{]} Halbherr, M., Zhou, Y., and Joerg, C. F. MIMD-style parallel
programming with continuation-passing threads. In Proceedings of the 2nd
International Workshop on Massive Parallelism: Hardware, Software, and
Applications, Capri, Italy, Sep.~1994.\\
{[}21{]} Halstead, Jr., R. H. Multilisp: A language for concurrent
symbolic computation. ACM Transactions on Programming Languages and
Systems, 7(4):501--538, Oct.~1985.\\
{[}22{]} Hillis, W. and Steele, G. Data parallel algorithms.
Communications of the ACM, 29(12):1170--1183, Dec.~1986.\\
{[}23{]} Hsieh, W. C., Wang, P., and Weihl, W. E. Computation migration:
Enhancing locality for distributed-memory parallel systems. In
Proceedings of the Fourth ACM SIGPLAN Symposium on Principles and
Practice of Parallel Programming (PPoPP), pp.~239--248, San Diego,
California, May 1993.\\
{[}24{]} Jagannathan, S. and Philbin, J. A customizable substrate for
concurrent languages. In Proceedings of the ACM SIGPLAN '92 Conference
on Programming Language Design and Implementation, pp.~55--67, San
Francisco, California, Jun.~1992.\\
{[}25{]} Joerg, C. and Kuszmaul, B. C. Massively parallel chess. In
Proceedings of the Third DIMACS Parallel Implementation Challenge,
Rutgers University, New Jersey, Oct.~1994. Available as
ftp://theory.lcs.mit.edu/pub/cilk/dimacs94.ps.Z.\\
{[}26{]} Kalé, L. V. The Chare kernel parallel programming system. In
Proceedings of the 1990 International Conference on Parallel Processing,
Volume II: Software, pp.~17--25, Aug.~1990.

{[}27{]} Karamcheti, V. and Chien, A. Concert---efficient runtime
support for concurrent object-oriented programming languages on stock
hardware. In Supercomputing '93, pp.~598--607, Portland, Oregon,
Nov.~1993.\\
{[}28{]} Karp, R. M. and Ramachandran, V. Parallel algorithms for
shared-memory machines. In van Leeuwen, J., (Ed.), Handbook of
Theoretical Computer Science---Volume A: Algorithms and Complexity,
chapter 17, pp.~869--941. MIT Press, Cambridge, Massachusetts, 1990.\\
{[}29{]} Karp, R. M. and Zhang, Y. Randomized parallel algorithms for
backtrack search and branch-and-bound computation. Journal of the ACM,
40(3):765--789, Jul.~1993.\\
{[}30{]} Kranz, D. A., Halstead, Jr., R. H., and Mohr, E. Mul-T: A
high-performance parallel Lisp. In Proceedings of the SIGPLAN '89
Conference on Programming Language Design and Implementation,
pp.~81--90, Portland, Oregon, Jun.~1989.\\
{[}31{]} Kuszmaul, B. C. Synchronized MIMD Computing. Ph.D.~thesis,
Department of Electrical Engineering and Computer Science, Massachusetts
Institute of Technology, May 1994. Available as MIT Laboratory for
Computer Science Technical Report MIT/LCS/TR-645 or
ftp://theory.lcs.mit.edu/pub/bradley/phd.ps.z.\\
{[}32{]} Leiserson, C. E., Abuhamdeh, Z. S., Douglas, D. C., Feynman, C.
R., Ganmukhi, M. N., Hill, J. V., Hillis, W. D., Kuszmaul, B. C.,
Pierre, M. A. S., Wells, D. S., Wong, M. C., Yang, S.-W., and Zak, R.
The network architecture of the Connection Machine CM-5. In Proceedings
of the Fourth Annual ACM Symposium on Parallel Algorithms and
Architectures, pp.~272--285, San Diego, California, Jun.~1992.\\
{[}33{]} Liu, P., Aiello, W., and Bhatt, S. An atomic model for
message-passing. In Proceedings of the Fifth Annual ACM Symposium on
Parallel Algorithms and Architectures, pp.~154--163, Velen, Germany,
Jun.~1993.\\
{[}34{]} Miller, R. C. A type-checking preprocessor for Cilk 2, a
multithreaded C language. Master's thesis, Department of Electrical
Engineering and Computer Science, Massachusetts Institute of Technology,
May 1995.\\
{[}35{]} Mohr, E., Kranz, D. A., and Halstead, Jr., R. H. Lazy task
creation: A technique for increasing the granularity of parallel
programs. IEEE Transactions on Parallel and Distributed Systems,
2(3):264--280, Jul.~1991.\\
{[}36{]} Nikhil, R. S. A multithreaded implementation of Id using P-RISC
graphs. In Proceedings of the Sixth Annual Workshop on Languages and
Compilers for Parallel Computing, number 768 in Lecture Notes in
Computer Science, pp.~390--405, Portland, Oregon, Aug.~1993.
Springer-Verlag.\\
{[}37{]} Nikhil, R. S. Cid: A parallel, shared-memory C for
distributed-memory machines. In Proceedings of the Seventh Annual
Workshop on Languages and Compilers for Parallel Computing, Aug.~1994.\\
{[}38{]} Pande, V. S., Joerg, C. F., Grosberg, A. Y., and Tanaka, T.
Enumerations of the hamiltonian walks on a cubic sublattice. Journal of
Physics A, 27, 1994.\\
{[}39{]} Rinard, M. C., Scales, D. J., and Lam, M. S. Jade: A
high-level, machine-independent language for parallel programming.
Computer, 26(6):28--38, Jun.~1993.\\
{[}40{]} Rudolph, L., Slivkin-Allalouf, M., and Upfal, E. A simple load
balancing scheme for task allocation in parallel machines. In
Proceedings of the Third Annual ACM Symposium on Parallel Algorithms and
Architectures, pp.~237--245, Hilton Head, South Carolina, Jul.~1991.

{[}41{]} Sunderam, V. S. PVM: A framework for parallel distributed
computing. Concurrency: Practice and Experience, 2(4):315--339,
Dec.~1990.\\
{[}42{]} Tanenbaum, A. S., Bal, H. E., and Kaashoek, M. F. Programming a
distributed system using shared objects. In Proceedings of the Second
International Symposium on High Performance Distributed Computing,
pp.~5--12, Spokane, Washington, Jul.~1993.\\
{[}43{]} Vandevoorde, M. T. and Roberts, E. S. WorkCrews: An abstraction
for controlling parallelism. International Journal of Parallel
Programming, 17(4):347--366, Aug.~1988.\\
{[}44{]} Wu, I.-C. and Kung, H. T. Communication complexity for parallel
divide-and-conquer. In Proceedings of the 32nd Annual Symposium on
Foundations of Computer Science, pp.~151--162, San Juan, Puerto Rico,
Oct.~1991.

\section{Short biographies of the
authors}\label{short-biographies-of-the-authors}

ROBERT (BOBBY) BLUMOFE received his Bachelor's degree from Brown
University in 1988 and his Ph.D.~from MIT in 1995. He started his
research career working on computer graphics with Andy van Dam at Brown,
and did his Ph.D.~work on algorithms and systems for parallel
multithreaded computing with Charles Leiserson at MIT. As part of this
dissertation work, Bobby developed an adaptive and fault tolerant
version of Cilk, called Cilk-NOW, that runs on networks of workstations.
Bobby is now an Assistant Professor at the University of Texas at
Austin, and he is continuing his work on Cilk and Cilk-NOW.

CHRISTOPHER F. JOERG received his B.S. and M.S. degrees in Computer
Science and Engineering from MIT in 1987 and 1990, respectively. He
expects to receive his Ph.D.~from MIT in January, 1996. His research
interests are in the areas of parallel systems and computer networks.

BRADLEY C. KUSZMAUL received two S.B. degrees in 1984, an S.M. degree in
1986, and a Ph.D.~degree in 1994, all from MIT. In 1987, midway through
his Ph.D.~program, he took a year off from MIT to serve as one of the
principal architects of the Connection Machine CM-5 at Thinking Machines
Corporation. In 1995, he joined the Departments of Computer Science and
Electrical Engineering at Yale University, where he is now Assistant
Professor. Prof.~Kuszmaul's work to solve systems problems in
high-performance computing spans a wide range of technology including
VLSI chips, interconnection networks, operating systems, compilers,
interpreters, algorithms, and applications.

CHARLES E. LEISERSON received the B.S. degree in computer science and
mathematics from Yale University, New Haven, Connecticut, in 1975 and
the Ph.D.~degree in computer science from Carnegie Mellon University,
Pittsburgh, Pennsylvania, in 1981. In 1981, he joined the faculty of the
Massachusetts Institute of Technology, Cambridge, Massachusetts. He is
now Professor of Computer Science and Engineering in the MIT Department
of Electrical Engineering and Computer Science and head of the
Supercomputing Technologies group in the MIT Laboratory for Computer
Science.

Prof.~Leiserson's research centers on the theory of parallel computing,
especially as it relates to engineering reality. Prof.~Leiserson
pioneered the development of VLSI theory and has written many papers on
VLSI algorithms, graph layout, and computer-aided design. His
contributions include the divide-and-conquer method of graph layout and
the retiming method for optimizing

digital circuitry. Prof.~Leiserson has been a leader in the development
of parallel computing. As a graduate student at Carnegie Mellon, he
wrote the first paper on systolic architectures with his advisor H.T.
Kung. While Corporate Fellow of Thinking Machines Corporation, he
designed and led the implementation of the network architecture for the
Connection Machine Model CM-5 Supercomputer, which incorporates the
fat-tree interconnection network he developed at MIT. He has designed
and engineered many parallel algorithms, including ones for matrix
linear algebra, graph algorithms, optimization, and sorting.
Prof.~Leiserson's recent work has focused on dynamic, asynchronous
parallel computing.

Prof.~Leiserson has won numerous awards. His Ph.D.~dissertation,
Area-Efficient VLSI Computation, which deals with the design of systolic
systems and with the problem of determining the VLSI area of a graph,
won the first ACM Doctoral Dissertation Award in 1981, as well as the
Fannie and John Hertz Foundation Doctoral Thesis Award. In 1985 he
received a Presidential Young Investigator Award from the National
Science Foundation. Three of his papers have received awards from the
IEEE International Conference on Parallel Processing. His textbook,
Introduction to Algorithms, coauthored with Ronald L. Rivest and Thomas
H. Cormen, was named Best 1990 Professional and Scholarly Book in
Computer Science and Data Processing by the Association of American
Publishers. Prof.~Leiserson is a member of the ACM, IEEE, and SIAM. He
currently serves as General Chair for the ACM Symposium on Parallel
Algorithms and Architectures. Prof.~Leiserson's proudest professional
accomplishments are the 18 students whose doctoral theses he has had the
privilege to supervise.

KEITH H. RANDALL received his B.S. and M.S. degrees in Computer Science
and Engineering from MIT in 1993 and 1994, respectively. He expects to
receive his Ph.D.~from MIT in 1997. His research interests include
routing, parallel algorithms, and scheduling.

YULI ZHOU received his B.S. in Electric Engineering from the University
of Science and Technology of China in 1983, and M.S. and Ph.D.~in
Computer Science from the Graduate School of Arts and Sciences, Harvard
University in 1990. Since then he has worked at the MIT Laboratory for
computer science as a research associate on compilers for parallel
programming languages in the Computation Structures Group. He is
currently a member of the technical staff at AT\&T Bell Laboratories.
His main research interests are programming languages and
parallel/distributed computing.

\end{document}
